\chapter*{ЗАКЛЮЧЕНИЕ}
\addcontentsline{toc}{chapter}{ЗАКЛЮЧЕНИЕ}
\label{chap:conclusion}

В работе показано, что линейный регрессионный анализ может использоваться для структурной диагностики нелинейных процессов. При этом область его содержательной применимости заметно уже области формальной подгонки. В рамках вычислительного эксперимента установлено, что регрессия надёжно восстанавливает структуру в устойчивых и информативных режимах, но начинает давать ложные лаги и терять интерпретируемость в циклических, переходных и геометрически вырожденных окнах \cite{strogatz,montgomery}.

Синтетическая часть позволила выделить пять принципиальных результатов. Базовая, запаздывающая и смешанная модели различаются по типу истинного фактора, и это различие регрессия распознаёт в устойчивых режимах. Высокий коэффициент детерминации не гарантирует правильной интерпретации структуры и тем более не гарантирует устойчивого итеративного прогноза. Короткое окно остаётся полезным локальным инструментом, но требует предварительной проверки режима, локальной дисперсии и лаговой геометрии. Дополнительные проверки устойчивости показывают, что ни регуляризация, ни PCA, ни настройка порога значимости не заменяют предварительную диагностику режима и интерпретируемости. В результате из синтетической части получен маршрут анализа короткого окна: сначала режим, затем интерпретируемость, далее выбор контура и только после этого чтение коэффициентов.

Эмпирическая часть подтвердила переносимость этого маршрута на реальные данные по выручке. На уровне групп фирм было показано, что короткие окна информативнее полных интервалов для содержательного чтения годовых денежных траекторий. На отраслевом уровне установлено, что сектора различаются и по качеству подгонки, и по допустимости анализа: часть отраслей допускает структурное чтение, другая часть должна рассматриваться только как фазовые или неинтерпретируемые окна. Главный эмпирический результат состоит в построении процедуры, которая заранее отделяет читаемые окна от тех, где коэффициенты будут вводить в заблуждение.

Практическая значимость исследования проявляется в трёх типовых применениях. Предложенный подход позволяет отбраковывать окна, где высокое качество подгонки создаёт лишь видимость устойчивой структуры. Он позволяет различать динамические режимы до запуска регрессии и заранее выбирать допустимый аналитический контур. Наконец, он делает аккуратнее анализ выручки групп и отраслей, где простое сравнение коэффициентов без предварительной диагностики режима может давать ложные выводы.

Ограничения исследования также были зафиксированы. Синтетические эксперименты проводились в контролируемых условиях и на ограниченном наборе моделей. На реальных данных точное восстановление истинной структуры невозможно, поэтому интерпретация строится на устойчивости режима, повторяемости паттернов и допустимости чтения окна. Кроме того, отраслевой анализ выручки ограничен неполнотой кодировки ОКВЭД и различием экономической природы денежных и промышленных рядов \cite{minec-okved,rosstat-okved,rosstat-ipp,hse-iipp}.

Поставленная цель достигнута. В работе показано, что линейный регрессионный анализ пригоден как инструмент структурной диагностики нелинейных процессов при условии предварительной режимной проверки и оценки интерпретируемости окна. В таком виде он может быть перенесён из синтетического эксперимента на реальные социально-экономические данные.
